\documentclass[11pt]{article}
\usepackage{amsmath,amssymb,amsthm}

\newtheorem{theorem}{Theorem}
\newtheorem{corollary}[theorem]{Corollary}
\newtheorem{lemma}[theorem]{Lemma}

\newcommand{\Up}{U_+}
\newcommand{\mass}{\operatorname{B}}
\newcommand{\car}{\operatorname{C}}
\newcommand{\sh}{\operatorname{D}}

\title{Sharp carry rearrangement for variable-depth inverse prefixes}
\author{}
\date{27 August 2026}

\begin{document}
\raggedbottom
\maketitle

For a positive odd integer $x$, put
\[
 \Up(x):=\frac{3x+1}{2^{v_2(3x+1)}}.
 \tag{1}
\]
For a positive exponent word $E=(e_1,\ldots,e_h)$, define
\[
 \mass(E):=\sum_{i=1}^h e_i,
 \qquad
 \car(E):=\sum_{i=1}^h3^{h-i}2^{e_1+\cdots+e_{i-1}},
 \tag{2}
\]
and
\[
 \sh(E):=\car(E)+2^{\mass(E)}-3^h.
 \tag{3}
\]
Thus an exact path with word $E$ from $z$ to $r$ satisfies
\[
 2^{\mass(E)}(r+1)=3^h(z+1)+\sh(E).
 \tag{4}
\]
Write
\[
 \mass(E)=h+t,
 \qquad t\geq0.
 \tag{5}
\]
The earlier variable-depth reduction used the bounds
$0\leq\sh(E)\leq3^h(2^t-1)$.  The purpose of this note is to replace
them with the exact extrema for fixed $h$ and $t$.

\section{The carry rearrangement}

\begin{lemma}[adjacent excess transfer]
Let $U,V$ be positive exponent words, either of which may be empty, and
let $a\geq2$ and $b\geq1$.  Then
\[
 \begin{aligned}
 &\sh\bigl(U(a,b)V\bigr)
   -\sh\bigl(U(a-1,b+1)V\bigr)\\
 &\hspace{35mm}
 =2^{\mass(U)+a-1}3^{|V|}>0.
 \end{aligned}
 \tag{6}
\]
In other words, moving one unit of excess one position to the right
strictly decreases the shifted carry.
\end{lemma}

\begin{proof}
For two letters,
\[
 \sh(a,b)=2^{a+b}+2^a-6.
 \tag{7}
\]
We also have the concatenation identity
\[
 \sh(XY)=3^{|Y|}\sh(X)+2^{\mass(X)}\sh(Y).
 \tag{8}
\]
Apply (8) twice to $U(a,b)V$.  Since $a+b$ is unchanged by the
transfer, every term except $2^{\mass(U)}3^{|V|}2^a$ cancels.  This
gives (6).
\end{proof}

\begin{theorem}[sharp shifted-carry extrema]
For every positive exponent word $E$ of length $h$ and excess $t$,
\[
 \boxed{
 2^h(2^t-1)\leq\sh(E)
 \leq2\,3^{h-1}(2^t-1).}
 \tag{9}
\]
If $t>0$ and $h>1$, equality on the left occurs exactly for
\[
 E=1^{h-1}(t+1),
 \tag{10}
\]
while equality on the right occurs exactly for
\[
 E=(t+1)1^{h-1}.
 \tag{11}
\]
For $t=0$, both bounds are equalities and $E=1^h$.
\end{theorem}

\begin{proof}
Write $e_i=1+c_i$, where the nonnegative integers $c_i$ have sum $t$.
If some $c_i>0$ with $i<h$, Lemma~1 moves one unit from $c_i$ to
$c_{i+1}$ and strictly decreases $\sh(E)$.  Repeating this operation
moves all excess to the final letter.  Therefore, the unique minimum
for $t>0$ and $h>1$ is (10).  Direct substitution gives
\[
 \sh\bigl(1^{h-1}(t+1)\bigr)=2^h(2^t-1).
 \tag{12}
\]

Applying the transfer in the opposite direction strictly increases the
shifted carry until all excess lies in the first letter.  This gives the
unique maximum (11), with
\[
 \sh\bigl((t+1)1^{h-1}\bigr)
 =2\,3^{h-1}(2^t-1).
 \tag{13}
\]
The cases $t=0$ and $h=1$ follow directly.
\end{proof}

\section{The sharp prefix interval}

For $h\geq1$, put
\[
 \mathcal A_h(r):=\left(\frac23\right)^h(r+1).
 \tag{14}
\]

\begin{corollary}[sharp all-excess prefix interval]
Let $E$ be an exact $h$-step word from a positive odd integer $z$ to a
positive odd integer $r$, and suppose that $\mass(E)=h+t$.  Then
\[
 \boxed{
 2^t\left(\mathcal A_h(r)-\frac23\right)-\frac13
 \leq z
 \leq
 \left(\frac23\right)^h(2^t r+1)-1.}
 \tag{15}
\]
The lower equality has word $(t+1)1^{h-1}$, and the upper equality has
word $1^{h-1}(t+1)$, whenever the corresponding word is exact.

Compared with the earlier interval, the lower endpoint is increased by
$(2^t-1)/3$, and the upper endpoint is decreased by
$(2/3)^h(2^t-1)$.  Both improvements are strict when $t>0$.
\end{corollary}

\begin{proof}
Equation (4) gives
\[
 z+1=2^t\mathcal A_h(r)-\frac{\sh(E)}{3^h}.
 \tag{16}
\]
Substitute the upper bound in (9) to obtain the lower bound in (15),
and substitute the lower bound in (9) to obtain the upper bound.  The
equality statements follow from Theorem~2.
\end{proof}

\begin{corollary}[linear excess floor at fixed terminal state]
Every exact positive $h$-step predecessor word of $r$ satisfies
\[
 \boxed{
 t\geq
 \max\left\{0,
 \left\lceil
 \log_2\left(\frac{2(3/2)^h-1}{r}\right)
 \right\rceil\right\}.}
 \tag{17}
\]
Consequently, along inverse paths of unbounded depth to one fixed
terminal state, the prefix excess tends to infinity with asymptotic
lower slope at least
\[
 \log_2\left(\frac32\right).
 \tag{18}
\]
\end{corollary}

\begin{proof}
Since $z$ is positive, the upper bound in (15) is at least one.  Hence
\[
 2^t r+1\geq2\left(\frac32\right)^h.
 \tag{19}
\]
Solving for the nonnegative integer $t$ proves (17), and (18) follows.
\end{proof}

\section{The remaining fixed-seed window}

We now apply the sharp interval to the variable-depth factorization.  Let
$q_*$ be the least same-depth suffix start, and let $r\geq q_*$ be the
start in any suffix channel.  If an $h$-step prefix of excess $t$ starts
at $z$, then (15) gives
\[
 z\geq
 2^t\left(\mathcal A_h(q_*)-\frac23\right)-\frac13.
 \tag{20}
\]

\begin{corollary}[sharp necessary condition for a smaller start]
Suppose that $\mathcal A_h(q_*)>2/3$.  If $z<N$, then
\[
 \boxed{
 2^t<
 \frac{N+1/3}{\mathcal A_h(q_*)-2/3}.}
 \tag{21}
\]
More generally, if $t\geq T$, then a start below $N$ is possible only if
\[
 \boxed{
 h>
 \log_{3/2}\left(
 \frac{q_*+1}{\displaystyle 2/3+(N+1/3)/2^T}
 \right).}
 \tag{22}
\]
For the first previously unclassified range, $T=3$, this becomes
\[
 \boxed{
 h>\log_{3/2}\left(\frac{24(q_*+1)}{3N+17}\right).}
 \tag{23}
\]
If a nonminimal suffix channel has the stronger lower bound $r\geq r_0$,
then $q_*$ may be replaced by $r_0$ in (20)--(23).
\end{corollary}

\begin{proof}
Inequality (21) is the contrapositive of (20).  If $t\geq T$, then
$2^t\geq2^T$, so (21) forces
\[
 \left(\frac23\right)^h(q_*+1)
 <\frac23+\frac{N+1/3}{2^T}.
 \tag{24}
\]
This is equivalent to (22).  Substitution of $T=3$ gives (23).
Monotonicity in the suffix start proves the last assertion.
\end{proof}

Let $n$ be the globally least positive odd integer with an unbounded
accelerated orbit.  Taking $N=n$ in Corollary~5 proves that no prefix of
excess at least three, in any suffix channel, can create a start below
$n$ before the explicit depth in (23).  Thus the unresolved branch is
strictly deeper than this threshold and has excess growing at least at
the linear rate in (17).  This does not yet exclude that deep branch.

\end{document}
